\documentclass[11pt]{article}

\usepackage[a4paper,margin=1in]{geometry}
\usepackage[T1]{fontenc}
\usepackage[utf8]{inputenc}
\usepackage{lmodern}
\usepackage{microtype}
\usepackage{amsmath,amssymb,amsthm,mathtools}
\usepackage{xurl}
\usepackage{authblk}
\usepackage{xcolor}
\usepackage[
  colorlinks=true,
  linkcolor=blue!55!black,
  citecolor=blue!55!black,
  urlcolor=blue!55!black,
  unicode=true,
  pdfauthor={Shouqiao Wang}
]{hyperref}
\hypersetup{pdftitle={A Proposed Solution to Erd\H{o}s Problem 788}}

\numberwithin{equation}{section}

\newtheorem{theorem}{Theorem}[section]
\newtheorem{proposition}[theorem]{Proposition}
\newtheorem{lemma}[theorem]{Lemma}
\newtheorem{corollary}[theorem]{Corollary}
\theoremstyle{definition}
\newtheorem{definition}[theorem]{Definition}
\theoremstyle{remark}
\newtheorem{remark}[theorem]{Remark}

\newcommand{\F}{\mathbb F}
\newcommand{\N}{\mathbb N}
\newcommand{\TV}{d_{\mathrm{TV}}}
\newcommand{\Ext}{\operatorname{Ext}}
\newcommand{\calC}{\mathcal C}

\title{A Proposed Solution to Erd\H{o}s Problem 788}
\author{Shouqiao Wang}
\affil{Columbia University\hspace{2.5em}Multiscalar Intelligence}
\date{}

\begin{document}
\maketitle

\begin{abstract}
For $n\ge1$, put
\[
 I_n=(n,2n)\cap\N,\qquad J_n=(2n,4n)\cap\N,
\]
let $f(n)$ be the largest integer $t$ such that every $B\subseteq J_n$
admits a set $C\subseteq I_n$ with $c+c'\notin B$ for distinct
$c,c'\in C$ and $|B|+|C|\ge t$.
We prove that there is an absolute constant $c>0$ such that, for all
sufficiently large $n$,
\[
 c\sqrt{n\log n}\le f(n)
 \le n^{\frac12+O\!\left(
   \left(\frac{\log\log n}{\log n}\right)^{1/3}\right)}.
\]
In particular, $f(n)=n^{1/2+o(1)}$, which gives an affirmative answer to
the principal question for every sufficiently large $n$.

The upper bound is obtained from a family of $p^{o(r)}$ surjective
$\F_p$-linear maps $\F_p^{2r}\to\F_p^r$ which is a strong seeded
extractor at min-entropy $r+o(r)$.  Taking the union of their kernels
gives a palette of $p^{r+o(r)}$ group sums with independence number
$p^{r+o(r)}$.  We give a self-contained mixed-alphabet Trevisan--RRV
reconstruction argument establishing the extractor, and then lift the
construction to ordinary integer sums while accounting for every base-$p$
carry.  An ordered block design uses the unused suffix of each hybrid to
pay for its overlap excess, which gives the stated quantitative rate.
The lower bound follows from the sparse-neighborhood coloring
theorem of Alon, Krivelevich, and Sudakov.
This proposed solution was found by GPT-5.
\end{abstract}

\section{Introduction}\label{sec:introduction}

Throughout, $\N=\{1,2,\ldots\}$, and all logarithms are natural unless a
base is explicitly indicated.  For $n\in\N$, put
\[
 I_n=(n,2n)\cap\N,
 \qquad
 J_n=(2n,4n)\cap\N.
\]
For $B\subseteq J_n$, call $C\subseteq I_n$ \emph{$B$-admissible} if
$c+c'\notin B$ whenever $c,c'\in C$ are distinct, and define
\[
 f(n)=\max\bigl\{t\in\mathbb Z:\ 
 \text{every $B\subseteq J_n$ has a $B$-admissible $C$ with
 $|B|+|C|\ge t$}\bigr\}.
\]

For a finite nonempty set $W$, write $U_W$ for its uniform distribution.
For probability distributions $P,Q$ on the same finite set, write
\[
 \TV(P,Q)=\frac12\sum_x|P(x)-Q(x)|.
\]
When $Q$ has full support, also write
\[
 D(P\Vert Q)=\sum_x P(x)\log\frac{P(x)}{Q(x)}.
\]
This is Problem 788 in the Erd\H{o}s Problems collection
\cite{ErdosProblems788}.  The collection cites Erd\H{o}s's 1973 survey
\cite{Erdos1973}, which credits Choi's 1971 paper \cite{Choi1971} with
raising the problem.

Choi proved $f(n)\ll n^{3/4}$ and conjectured the bound
$f(n)\ll_\varepsilon n^{1/2+\varepsilon}$ \cite{Choi1971}.  Baltz, Schoen, and Srivastav
improved the upper bound to $f(n)\ll(n\log n)^{2/3}$
\cite{BSS2000}.  The online problem record reports the subsequent bound
$f(n)\le n^{3/5+o(1)}$, obtained by combining a random-Cayley-graph
reduction with work of Alon and Pham
\cite{ErdosProblems788,AlonPham2025}.  None of these
upper bounds is used below; our construction directly gives the conjectured
exponent for every sufficiently large interval length.

Our main theorem gives a little more quantitative information than the
requested exponent.

\begin{theorem}[Main theorem]\label{thm:main}
There are absolute constants $c,C>0$ such that, for all sufficiently
large $n$,
\[
 c\sqrt{n\log n}\le f(n)
 \le n^{\frac12+C\left(
   \frac{\log\log n}{\log n}\right)^{1/3}}.
\]
Consequently
\[
 f(n)=n^{1/2+o(1)}.
\]
\end{theorem}

\begin{remark}[Quantifiers in the upper bound]\label{rem:quantifiers}
The theorem supplies a witness palette for every sufficiently large
integer $n$, not merely along a subsequence.  Equivalently, for each fixed
$\varepsilon>0$ and all sufficiently large $n$, one can choose a single
$B\subseteq J_n$ such that every $B$-admissible $C\subseteq I_n$ obeys
$|B|+|C|\le n^{1/2+\varepsilon}$.
\end{remark}

The two sides of the proof use different information.  The lower bound
holds for every palette $B$.  For the upper bound we construct one palette,
depending on $n$, by combining coding theory, weak designs, and an exact
carry closure.

The proof is organized as follows.  Section~\ref{sec:graph} converts the
problem into an exact minimization involving the independence number of a
sum graph.  Section~\ref{sec:lower} proves the universal lower bound.
Section~\ref{sec:extractor} constructs and verifies a mixed-alphabet linear
seeded extractor, using an ordered weak design whose overlap excess is paid
for by the remaining hybrid coordinates;
Section~\ref{sec:group-palette} converts its kernels into a finite-field sum
palette.  Section~\ref{sec:carries} lifts that palette to ordinary integer
sums without discarding any carry pattern.
Section~\ref{sec:all-N} chooses the parameters uniformly for every interval
length, and Section~\ref{sec:completion} translates the construction back
to the original intervals and completes the proof.

\section{Exact graph formulation}\label{sec:graph}

For $s\in J_n$, let $M_s$ be the graph on $I_n$ whose edges are the
unordered pairs $\{c,c'\}$ satisfying
\[
 c\ne c',\qquad c+c'=s.
\]
Each $M_s$ is a matching: after one endpoint is specified, the other is
forced.  A possible equality $c=s-c$ creates no loop, since the problem
only constrains distinct elements.  For $B\subseteq J_n$, put
\[
 G_B=\bigcup_{s\in B}M_s.
\]

\begin{proposition}[Min--max identity]\label{prop:minmax}
For every $n\ge1$,
\[
 \boxed{\displaystyle
 f(n)=\min_{B\subseteq J_n}\bigl(|B|+\alpha(G_B)\bigr).}
\]
Moreover, the minimum is unchanged if $B$ is required to consist only
of sums of two distinct elements of $I_n$.
\end{proposition}

\begin{proof}
A set $C\subseteq I_n$ is $B$-admissible exactly when it is independent
in $G_B$.  Thus, for fixed $B$, the largest possible value of
$|B|+|C|$ is $|B|+\alpha(G_B)$.  An integer $t$ has the universal
property in the definition of $f(n)$ exactly when
\[
 t\le |B|+\alpha(G_B)\qquad\text{for every }B\subseteq J_n.
\]
The family of choices of $B$ is finite, so taking the largest such $t$
gives the displayed minimum.

If $s$ is not the sum of two distinct elements of $I_n$, then $M_s$ is
empty.  Deleting such an $s$ from $B$ leaves $G_B$ unchanged and decreases
$|B|$.  Hence no minimizer needs an unattainable sum.
\end{proof}

It will be useful to normalize the vertices.  Put
\[
 N=n-1,\qquad [N]_0=\{0,1,\ldots,N-1\}.
\]
Write $c=n+1+x$.  Then
\[
 c+c'=2n+2+x+y.
\]
For $N\ge2$, the sums of two distinct members of $[N]_0$ are exactly
\[
 \{1,2,\ldots,2N-3\}.
\]
For $A\subseteq\{1,\ldots,2N-3\}$, let $H_A$ be the graph on $[N]_0$
in which distinct $x,y$ are adjacent exactly when $x+y\in A$.  Translation
identifies $H_A$ with $G_B$ for
\[
 B=2n+2+A\subseteq\{2n+3,\ldots,4n-3\}\subseteq J_n.
\]
We will construct and estimate $H_A$ and translate back only at the end.
If $b=|A|$, then
\begin{equation}\label{eq:max-degree}
 \Delta(H_A)\le b,
\end{equation}
because each selected sum supplies at most one neighbor of any fixed
vertex.

\section{Universal lower bounds}\label{sec:lower}

We first record an elementary bound which makes the square-root scale
visible.

\begin{lemma}[Edge sums control chromatic number]\label{lem:chromatic}
Let a simple graph have distinct real labels on its vertices, and suppose
that its edges use $b$ distinct sums of endpoint labels.  Then
\[
 \chi(G)\le \left\lfloor\frac{b+3}{2}\right\rfloor.
\]
\end{lemma}

\begin{proof}
Put $k=\chi(G)$ and pass to a vertex-minimal induced subgraph $K$ with
chromatic number $k$.  Then $\delta(K)\ge k-1$.  Let $u$ and $v$ be the
least and greatest labels in $K$.  The at least $k-1$ sums on edges
incident with $u$ are distinct and at most $u+v$.  The at least $k-1$
sums on edges incident with $v$ are distinct and at least $u+v$.  The
two collections intersect in at most the one value $u+v$.  Hence
$b\ge2k-3$, which is equivalent to the assertion.
\end{proof}

Taking the largest color class in $H_A$, where $b=|A|$, gives
\[
 \alpha(H_A)\ge
 \left\lceil\frac{N}{\lfloor(b+3)/2\rfloor}\right\rceil.
\]
Optimizing this elementary expression yields, for $N\ge4$,
\[
 b+\alpha(H_A)\ge \lceil\sqrt{8N}\rceil-3.
\]
The case $b=0$ gives the value $N$ and cannot minimize when $N\ge4$
(the choice $A=\{N-1\}$ has $b=1$ and gives
$1+\lceil N/2\rceil\le N$).  For $b\ge1$ the
relevant values are $b=2q-3$, and
\[
 \min_{q\ge2}\left(2q-3+\left\lceil\frac Nq\right\rceil\right)
 =\lceil\sqrt{8N}\rceil-3.
\]
For completeness, put $t=\lceil\sqrt{8N}\rceil$.  The inequality
$2q+N/q\ge\sqrt{8N}$ gives the lower direction.  Conversely,
\[
 \max_{q\in\mathbb Z}q(t-2q)=\left\lfloor\frac{t^2}{8}\right\rfloor
 \ge N.
\]
For $N\ge4$ a maximizing positive integer $q$ can be chosen with $q\ge2$;
then $N/q\le t-2q$ and hence
$2q+\lceil N/q\rceil\le t$.  This proves the identity.
This already proves $f(n)\ge2\sqrt{2n}-O(1)$.  We next gain a logarithmic
factor.

\begin{lemma}[Triangle sum triples]\label{lem:triangles}
If $H_A$ uses $b$ selected sums and has $T$ triangles, then
\[
 T\le\binom b3.
\]
\end{lemma}

\begin{proof}
A triangle with vertices $x,y,z$ has the three pairwise distinct edge
sums
\[
 x+y,\qquad x+z,\qquad y+z.
\]
Their unordered set determines the triangle: if the three incident sums
are denoted by $s_{xy},s_{xz},s_{yz}$, then
\[
 x=\frac{s_{xy}+s_{xz}-s_{yz}}2,
 \quad
 y=\frac{s_{xy}+s_{yz}-s_{xz}}2,
 \quad
 z=\frac{s_{xz}+s_{yz}-s_{xy}}2.
\]
Equivalently, from an unordered triple $\{a,b,c\}$ one recovers the
unordered vertex set
\[
 \left\{\frac{a+b-c}{2},\frac{a+c-b}{2},\frac{b+c-a}{2}\right\}.
\]
Thus triangles inject into the three-element subsets of $A$.
\end{proof}

We use the following conservative form of a theorem of Alon,
Krivelevich, and Sudakov
\cite[Theorem~1.1]{AKS1999}.

\begin{theorem}[Sparse-neighborhood coloring theorem, safe form]\label{thm:aks}
There is an absolute $K>0$ with the following property.  If a graph has
maximum degree at most $d$ and every vertex neighborhood spans at most
$d^2/F$ edges, where $2\le F\le d^2$, then
\[
 \chi(G)\le K\frac d{\log F}.
\]
\end{theorem}

The closed endpoint $F=d^2$ follows from the usual strict-range formulation
by applying it with $F'=d^2/2$ and changing the absolute constant; bounded
$d$ is absorbed by greedy coloring.
The primary statement calls $d$ the maximum degree.  If that convention
is read as requiring equality rather than an upper bound, one may adjoin
a disjoint star $K_{1,d}$: this makes the maximum degree exactly $d$,
adds no edges to any neighborhood, and does not affect the coloring bound
for the original component.

\begin{proposition}[Universal lower bound]\label{prop:loglower}
There is an absolute $c>0$ such that, for every sufficiently large $N$
and every $A\subseteq\{1,\ldots,2N-3\}$,
\[
 |A|+\alpha(H_A)\ge c\sqrt{N\log N}.
\]
Consequently $f(n)\gg\sqrt{n\log n}$.
\end{proposition}

\begin{proof}
Write $b=|A|$.  For a vertex $v$, let $t_v$ be the number of edges
spanned by its neighborhood, equivalently the number of triangles
containing $v$.  Lemma~\ref{lem:triangles} gives
\[
 \sum_v t_v=3T\le3\binom b3\le\frac{b^3}{2}.
\]
Therefore the set
\[
 U=\left\{v:t_v\le\frac{b^3}{N}\right\}
\]
has at least $N/2$ vertices.  The induced graph $H=H_A[U]$ has maximum
degree at most $b$, and every neighborhood in $H$ spans at most $b^3/N$
edges.

Suppose $2\le b\le N/2$, and put
\[
 F_0=\min\left\{b^2,\frac Nb\right\}\ge2.
\]
In either case $b^3/N\le b^2/F_0$.  Applying
Theorem~\ref{thm:aks} to $H$ and taking a largest color class gives
\begin{equation}\label{eq:alphalower}
 \alpha(H_A)\ge
 c_0\frac Nb\log\!\left(\min\left\{b^2,\frac Nb\right\}\right)
\end{equation}
for an absolute $c_0>0$.

Now split into three ranges.  If $b\ge\sqrt{N\log N}$, the palette
itself proves the result.  If $b<N^{1/4}$, the maximum-degree bound
\eqref{eq:max-degree}, followed by greedy coloring, gives
\[
 \alpha(H_A)\ge\frac{N}{b+1}\gg\sqrt{N\log N}.
\]
In the remaining range
\[
 N^{1/4}\le b<\sqrt{N\log N},
\]
we have
\[
 \min\left\{b^2,\frac Nb\right\}
 \ge\sqrt{\frac{N}{\log N}},
\]
and hence its logarithm is at least $(1/3)\log N$ for large $N$.
Equation~\eqref{eq:alphalower} now yields
\[
 b+\alpha(H_A)\ge b+c_1\frac{N\log N}{b}
 \ge2\sqrt{c_1N\log N}.
\]
The cases $b=0,1$ are covered by the maximum-degree estimate.  The claim
for $f$ follows from Proposition~\ref{prop:minmax}.
\end{proof}

\section{A linear seeded extractor over a growing odd field}
\label{sec:extractor}

The upper bound rests on the following finite-field statement.  All
logarithms used in min-entropy in this section are to base $p$.

\begin{theorem}[Linear extractor family]\label{thm:linear-family}
Let $p$ run through odd primes and let $r\to\infty$ in such a way that
\[
 \bigl(\log(pr)\bigr)^2\log\bigl(2+\log(pr)\bigr)
 =o(r\log p).
\]
Then there are integers $d,s=o(r)$ and an indexed family
\[
 \{F_y:\F_p^{2r}\to\F_p^r:y\in\mathcal Y\},
 \qquad |\mathcal Y|\le p^d,
\]
such that every $F_y$ is surjective and $\F_p$-linear, and every random
variable $X$ on $\F_p^{2r}$ with
\[
 H_{\infty,p}(X):=-\log_p\max_x\Pr[X=x]\ge r+s
\]
satisfies
\[
 \frac1{|\mathcal Y|}\sum_{y\in\mathcal Y}
 \TV(F_y(X),U_{\F_p^r})<\frac13.
\]
More quantitatively, there is an auxiliary integer
$\ell=O(\log(pr))$ for which one may take
\[
 d,s=O\left(
 \frac{\ell\sqrt r+\ell^2\log(\ell+1)}{\log p}
 +\log_p(pr)+1\right).
\]
The family is indexed: repeated maps are allowed.
\end{theorem}

The construction is a mixed-alphabet adaptation of the linear form of
Trevisan's extractor \cite{Trevisan2001} and the weak-design
reconstruction analysis of Raz, Reingold, and Vadhan \cite{RRV2002}.
The code and every fixed-seed output map are $\F_p$-linear, but code
coordinates and seed coordinates are indexed by binary strings.  This
decoupling is useful: an overlap of $t$ seed bits has $2^t$ possible inputs,
independently of the growing output alphabet.  We prove the required code,
design, and reconstruction statements with this distinction explicit.

The proof proceeds in four stages.  We first construct a short linear code,
and then a binary weak design controlling its overlap moments.  A hybrid
argument and reconstruction count combine these objects into a strong linear
extractor.  Finally we choose the parameters and discard rank-deficient seeds;
Section~\ref{sec:group-palette} turns the retained maps into a sum palette by
taking the union of their kernels.

\subsection{A short linear list-decodable code}

We identify a word of length $2^\ell$ with a function
$\{0,1\}^\ell\to\F_p$.  Agreement always means the fraction of coordinates
on which two words have the same symbol.

\begin{lemma}[Short binary-coordinate linear code]\label{lem:code}
Let $p$ be a prime, $m\ge1$, and $0<\eta<1/2$.  There exist a positive integer
$\ell$ and an injective $\F_p$-linear map
\[
 \mathcal E:\F_p^m\longrightarrow
 \F_p^{\{0,1\}^\ell}
\]
such that, for every word $Q:\{0,1\}^\ell\to\F_p$,
\[
 \#\left\{x\in\F_p^m:
 \Pr_{z\sim U_{\{0,1\}^\ell}}[\mathcal E(x)(z)=Q(z)]
 >\frac1p+\eta\right\}
 \le J:=\frac2{\eta^2}+1.
\]
One can arrange
\begin{equation}\label{eq:code-length}
 2^\ell\le \frac{C_0m\log p}{\eta^4}
\end{equation}
for an absolute constant $C_0$.
\end{lemma}

\begin{proof}
Set
\[
 \tau=\frac{p\eta^2}{2(p-1)},
 \qquad \Delta=1-\frac1p-\tau.
\]
We first find a linear code of relative distance at least $\Delta$.  Write
\[
 H_p(a)=a\log_p(p-1)-a\log_p a-(1-a)\log_p(1-a)
\]
for the $p$-ary entropy function, with the usual continuous extension at
the endpoints.  Then
\[
 1-H_p(\Delta)=\frac{D(P\Vert U_{\F_p})}{\log p},
\]
where
\[
 P=\left(\frac1p+\tau,
 \frac1p-\frac{\tau}{p-1},\ldots,
 \frac1p-\frac{\tau}{p-1}\right).
\]
The total variation distance between $P$ and $U_{\F_p}$ is $\tau$.  Pinsker's
inequality, with natural logarithms \cite{CoverThomas2006}, gives
\[
 1-H_p(\Delta)\ge\frac{2\tau^2}{\log p}.
\]

Let $R=2^\ell$ be the least power of $2$ strictly exceeding
\[
 \frac{(m+1)\log p}{2\tau^2}.
\]
The displayed quantity exceeds $1$, so $\ell\ge1$.
Choose an $R\times m$ matrix $T$ uniformly over $\F_p$.  For every fixed
nonzero $x$, the word $Tx$ is uniform on $\F_p^R$.  The standard $p$-ary
Hamming-ball estimate \cite{CoverThomas2006} is exact in the form
\[
 \sum_{0\le j\le aR}\binom Rj(p-1)^j\le p^{RH_p(a)}
 \qquad\left(0<a\le1-\frac1p\right).
\]
Indeed, with $z=a/((p-1)(1-a))\le1$, the left side is at most
$z^{-aR}(1+(p-1)z)^R=p^{RH_p(a)}$.  Taking $a=\Delta$ and using
the preceding Pinsker bound show that
\[
 \Pr\bigl[\operatorname{wt}(Tx)<\Delta R\bigr]
 \le p^{-R(1-H_p(\Delta))}.
\]
A union bound over fewer than $p^m$ nonzero $x$ is strictly below one.
We may therefore fix $T$ for which every nonzero codeword has relative
weight at least $\Delta$.  This also makes $T$ injective.  Index the $R$
coordinates by $\{0,1\}^\ell$ and call the resulting encoder $\mathcal E$.
The choice of $R$, the formula for $\tau$, and rounding to the next power
of $2$ give, explicitly,
\[
 R\le 2\frac{(m+1)\log p}{2\tau^2}
 \le\frac{8m\log p}{\eta^4},
\]
which proves \eqref{eq:code-length}.

It remains to prove the list bound.  Represent the $p$ symbols by unit
vectors $v_a$ forming a regular simplex:
\[
 \langle v_a,v_b\rangle=
 \begin{cases}
 1,&a=b,\\
 -1/(p-1),&a\ne b.
 \end{cases}
\]
Represent a word by the normalized concatenation of its simplex vectors.
Two words with agreement fraction $a$ then have inner product
\[
 \frac{pa-1}{p-1}.
\]
Suppose $L$ codewords agree with $Q$ on more than $1/p+\eta$ of the
coordinates.  Let their unit vectors be $u_1,\ldots,u_L$, and let $w$ be
the vector of $Q$.  By this identity,
\[
 \langle u_i,w\rangle>\gamma:=\frac{p\eta}{p-1}.
\]
The minimum distance gives, for $i\ne j$,
\[
 \langle u_i,u_j\rangle\le
 \beta:=\frac{p\tau}{p-1}=\frac{\gamma^2}{2}.
\]
Cauchy--Schwarz now implies
\[
 L^2\gamma^2
 <\left\langle\sum_i u_i,w\right\rangle^2
 \le\left\|\sum_i u_i\right\|^2
 \le L+L(L-1)\beta.
\]
Thus
\[
 L\le\frac{1-\beta}{\gamma^2-\beta}
 \le\frac2{\gamma^2}\le\frac2{\eta^2},
\]
as required.
\end{proof}

\subsection{An ordered binary weak design}

\begin{lemma}[Suffix-slack design]\label{lem:design}
For all integers $\ell,r\ge1$, there are an integer $D$ and
$\ell$-subsets $S_1,\ldots,S_r\subseteq[D]$ such that, for every $i$,
\begin{equation}\label{eq:weak-design}
 e_i:=\sum_{j<i}\left(2^{|S_i\cap S_j|}-1\right)\le r-i,
\end{equation}
and
\[
 D=O\left(\ell\sqrt r+\ell^2\log(\ell+1)\right).
\]
The implied constant is absolute, and the sets and their ordering may be
chosen deterministically from $(r,\ell)$.
\end{lemma}

\begin{proof}
Partition the $r$ rows into consecutive blocks.  If $R_b$ rows remain
before block $b$, put
\[
 m_b=\left\lceil\frac{R_b}{2}\right\rceil,
 \qquad
 R_{b+1}=\left\lfloor\frac{R_b}{2}\right\rfloor.
\]
Thus $R_{b+1}\ge m_b-1$.  Give different blocks disjoint seed-coordinate
universes.

For a block of size $m=m_b$, let $q$ be the least prime at least
\[
 Q=\max\{2,\ell,\lceil\sqrt m\rceil\}.
\]
Bertrand's postulate gives $q<2Q$ \cite{HardyWright2008}.  Choose an
$\ell$-set $X\subseteq\F_q$ and use the universe $X\times\F_q$.
There are $q^2\ge m$ affine functions $x\mapsto ax+b$; select any $m$
distinct ones and associate to $(a,b)$ its graph
\[
 S_{a,b}=\{\,(x,ax+b):x\in X\,\}.
\]
Two distinct affine graphs meet in at most one point.

Let a row have zero-based local index $t$ in its block and one-based global
index $i$.  Rows in earlier blocks have disjoint universes and contribute
zero to the excess in \eqref{eq:weak-design}.  Each of the $t$ earlier
block-mates contributes at most $2^1-1=1$.  Moreover, the exact number of
global rows after this row is
\[
 r-i=(m_b-t-1)+R_{b+1}.
\]
Consequently
\[
 e_i\le t\le m_b-1\le R_{b+1}\le r-i.
\]
This includes the final global row, for which the last block has size one
and the excess is zero.

The universe of block $b$ has size
\[
 \ell q=O\bigl(\ell\max\{\ell,\sqrt{m_b}\}\bigr).
\]
The block sizes decrease geometrically, so
$\sum_b\sqrt{m_b}=O(\sqrt r)$.  Once $m_b<\ell^2$, only
$O(\log(\ell+1))$ nonempty blocks remain, each using $O(\ell^2)$
coordinates.  The disjoint union of the block universes therefore has the
claimed size; identify it with $[D]$.  Taking the least eligible prime and
the first affine functions in lexicographic order makes the construction
deterministic.
\end{proof}

The subtracted $1$ in \eqref{eq:weak-design} isolates the unavoidable
baseline: even disjoint sets contribute $2^0=1$ to the reconstruction
count.  Equivalently,
\[
 \sum_{j<i}2^{|S_i\cap S_j|}\le r-1.
\]
The remaining $r-i$ output coordinates will pay for the overlap excess at
hybrid row $i$.

\subsection{The reconstruction argument}

Fix
\begin{equation}\label{eq:eps-h-eta}
 \varepsilon=\frac1{20},\qquad
 h=\frac{\varepsilon}{r(p-1)},\qquad
 \eta=\frac h2.
\end{equation}
Apply Lemma~\ref{lem:code} with input dimension $2r$, and let
$\mathcal E$ be the resulting encoder, with coordinate set $\{0,1\}^\ell$
and list size $J=2/\eta^2+1$.  Apply Lemma~\ref{lem:design} with these
values of $\ell$ and $r$.
Put
\[
 d=\left\lceil D\log_p2\right\rceil,
\]
so that the binary seed set has size $2^D\le p^d$.  For
$x\in\F_p^{2r}$ and $y\in\{0,1\}^D$, define
\begin{equation}\label{eq:ext-definition}
 \Ext(x,y)_i=\mathcal E(x)(y|_{S_i}),
 \qquad 1\le i\le r.
\end{equation}
The coordinates of each $S_i$ are read in increasing order.

\begin{proposition}[Extractor inequality]\label{prop:extractor}
Suppose the integer $s$ satisfies
\begin{equation}\label{eq:s-threshold}
 s\ge d+\left\lceil\log_p\frac{4J}{h}\right\rceil+2,
 \qquad r+s\le2r.
\end{equation}
If $X$ is any random variable on $\F_p^{2r}$ with
$H_{\infty,p}(X)\ge r+s$, and $Y$ is uniform on $\{0,1\}^D$ independently
of $X$, then
\begin{equation}\label{eq:strong-ext}
 \TV\bigl((Y,\Ext(X,Y)),(Y,U_{\F_p^r})\bigr)\le\varepsilon.
\end{equation}
Here the copy of $U_{\F_p^r}$ on the right is independent of $Y$.
For each fixed seed $y$, the map $x\mapsto\Ext(x,y)$ is
$\F_p$-linear with codomain $\F_p^r$.
\end{proposition}

\begin{proof}
Write $Z=\Ext(X,Y)=(Z_1,\ldots,Z_r)$ and suppose that
\eqref{eq:strong-ext} fails.  Define the hybrids
\[
 \mathsf H_i=(Y,Z_1,\ldots,Z_i,U_{i+1},\ldots,U_r),
 \qquad 0\le i\le r,
\]
where the displayed $U_j$ are independent uniform symbols, also
independent of $(X,Y,Z)$.  The triangle inequality gives
an $i$ such that
\[
 \mathbb E_{(Y,Z_{<i})}
 \TV\bigl(\mathcal L(Z_i\mid Y,Z_{<i}),U_{\F_p}\bigr)
 >\frac{\varepsilon}{r}.
\]
For a distribution $\mu$ on $\F_p$,
\[
 \max_a\mu(a)-\frac1p
 \ge\frac{\TV(\mu,U_{\F_p})}{p-1},
\]
because the total positive deviation is the total variation distance and
has at most $p-1$ positive summands.  At every context choose a most
likely symbol.  The preceding inequality produces a deterministic
predictor
\[
 \mathcal P:\{0,1\}^D\times\F_p^{i-1}\to\F_p
\]
such that
\[
 \Pr[\mathcal P(Y,Z_{<i})=Z_i]>\frac1p+h.
\]
The index $i$ and the predictor $\mathcal P$ may depend on the fixed law of
$X$, but from now on they are fixed and do not depend on the point $x$.

For fixed $x$, let $q_x$ be the success probability in the preceding
predictor bound over $Y$, with $Z=\Ext(x,Y)$, and put
\[
 \mathcal B=\left\{x:q_x>\frac1p+\frac h2\right\}.
\]
Since $q_x\le1$, averaging that predictor bound gives
\begin{equation}\label{eq:bad-mass-lower}
 \Pr[X\in\mathcal B]>\frac h2.
\end{equation}
Indeed, if $\theta=\Pr[X\in\mathcal B]$, then
\[
 \frac1p+h<\mathbb E q_X
 \le\frac1p+\frac h2+
 \theta\left(1-\frac1p-\frac h2\right),
\]
which implies $\theta>h/2$.

We now count $\mathcal B$.  Fix $x\in\mathcal B$.  Regard a binary assignment
on $S_i$ as a code coordinate in $\{0,1\}^\ell$, using the increasing order
on $S_i$.  By averaging over the seed coordinates outside $S_i$, we obtain a fixing
$a\in\{0,1\}^{[D]\setminus S_i}$ for which the predictor agrees with the
word $\mathcal E(x)$ on more than a $1/p+h/2=1/p+\eta$ fraction of
assignments $z\in\{0,1\}^{S_i}$.

After this outside fixing, the earlier output
$\mathcal E(x)(y|_{S_j})$ is a function of only the
$|S_i\cap S_j|$ overlap coordinates of $z$.  We may enumerate every such
function: a function $\{0,1\}^t\to\F_p$ has $p^{2^t}$ possibilities.  Thus
the number of descriptions consisting of the outside fixing and all
earlier-output functions is at most
\begin{align*}
 2^{D-\ell}\prod_{j<i}p^{2^{|S_i\cap S_j|}}
 &\le p^{d+\sum_{j<i}2^{|S_i\cap S_j|}}\\
 &=p^{d+(i-1)+e_i}
 \le p^{d+r-1}.
\end{align*}
For a description $\mathcal D=(a,(g_j)_{j<i})$ and an assignment
$z\in\{0,1\}^{S_i}$, let $y(a,z)$ be the unique seed whose restrictions to
$S_i$ and $[D]\setminus S_i$ are respectively $z$ and $a$, and define
\[
 Q_{\mathcal D}(z)=\mathcal P\bigl(
 y(a,z),(g_j(z|_{S_i\cap S_j}))_{j<i}\bigr).
\]
Via the increasing-order identification
$\{0,1\}^{S_i}\cong\{0,1\}^\ell$, this is one received word
$Q_{\mathcal D}:\{0,1\}^\ell\to\F_p$.  For each
$x\in\mathcal B$, the description obtained above gives a word $Q_{\mathcal D}$
which agrees with $\mathcal E(x)$ on more than a $1/p+\eta$ fraction of
coordinates.  Thus $x$ belongs to the agreement list of at least one
counted description.  Lemma~\ref{lem:code} bounds every such list by
$J$, so
\[
 |\mathcal B|\le Jp^{d+r-1}.
\]
The min-entropy assumption now gives
\begin{align*}
 \Pr[X\in\mathcal B]
 &\le p^{-r-s}|\mathcal B|\\
 &\le Jp^{d-s-1}<\frac h4,
\end{align*}
where \eqref{eq:s-threshold} gives explicitly
\[
 Jp^{d-s-1}
 \le Jp^{-\lceil\log_p(4J/h)\rceil-3}
 \le\frac{h}{4p^3}<\frac h4.
\]
This
contradicts \eqref{eq:bad-mass-lower}, proving \eqref{eq:strong-ext}.

Finally, for fixed $y$, every coordinate in \eqref{eq:ext-definition} is
one coordinate of the $\F_p$-linear encoding $\mathcal E$.  The tuple of
the $r$ selected coordinates is therefore $\F_p$-linear.  Surjectivity is
not asserted yet; it is enforced below.
\end{proof}

\subsection{Parameters and surjectivity}

From Lemma~\ref{lem:code}, \eqref{eq:eps-h-eta}, and the input dimension
$2r$,
\[
 2^\ell=O\bigl(r\log p\,(pr)^4\bigr).
\]
Therefore
\begin{equation}\label{eq:ell-bound}
 \ell=O(\log(pr)).
\end{equation}
Lemma~\ref{lem:design} gives
\[
 D=O\left(\ell\sqrt r+\ell^2\log(\ell+1)\right)
\]
and therefore
\[
 d=O\left(
 \frac{\ell\sqrt r+\ell^2\log(\ell+1)}{\log p}+1\right).
\]
Also
\[
 \log_p(4J/h)=O(1+\log_p(pr)),
\]
so \eqref{eq:s-threshold} holds with
\[
 s=O\left(
 \frac{\ell\sqrt r+\ell^2\log(\ell+1)}{\log p}
 +\log_p(pr)+1\right).
\]
The growth hypothesis in Theorem~\ref{thm:linear-family}, together with
\eqref{eq:ell-bound}, makes both $d$ and $s$ equal to $o(r)$.
Indeed, if $u=\log(pr)$ and $\lambda=\log p$, then
\[
 \left(\frac{u}{\lambda\sqrt r}\right)^2
 =\frac{u^2}{r\lambda^2}=o(1),
 \qquad
 \frac{u^2\log(2+u)}{r\lambda}=o(1),
\]
while $\log_p(pr)/r=o(1)$.  These three estimates control respectively
the $\ell\sqrt r$, $\ell^2\log(\ell+1)$, and remaining terms above.
For the first estimate, use the second one together with
$\lambda\ge\log3$ and $\log(2+u)>1$.
This verifies all parameter claims in
Theorem~\ref{thm:linear-family}, except surjectivity.

The full set of $2^D$ binary seeds is an indexed family of linear maps.  Apply
Proposition~\ref{prop:extractor} to the source $U_{\F_p^{2r}}$.
For any source $X$, the joint distance in \eqref{eq:strong-ext} is exactly
the fixed-seed average
\[
 \TV\bigl((Y,\Ext(X,Y)),(Y,U_{\F_p^r})\bigr)
 =\frac1{2^D}\sum_{y\in\{0,1\}^D}
 \TV(\Ext(X,y),U_{\F_p^r}).
\]
If a seed map has rank $t<r$, its output is uniform on a $p^t$-element
subspace, and therefore
\[
 \TV(\Ext(U_{\F_p^{2r}},y),U_{\F_p^r})
 =1-p^{t-r}\ge1-\frac1p.
\]
If $b$ is the fraction of rank-deficient seeds, \eqref{eq:strong-ext}
gives
\[
 b\le\frac{\varepsilon}{1-1/p}\le\frac3{40}.
\]
Retain the good seeds, with their original multiplicities, and put the
uniform distribution on this retained index set $\mathcal Y$.  For each
$y\in\mathcal Y$, set
\[
 F_y(x)=\Ext(x,y).
\]
Then $|\mathcal Y|\le p^d$, every $F_y$ is surjective, and for every
source covered by Proposition~\ref{prop:extractor},
\[
 \frac1{|\mathcal Y|}\sum_{y\in\mathcal Y}
 \TV(F_y(X),U_{\F_p^r})
 \le\frac{\varepsilon}{37/40}=\frac2{37}<\frac13.
\]
This completes the proof of Theorem~\ref{thm:linear-family}.

\section{From the extractor to a finite-field sum palette}
\label{sec:group-palette}

Put
\[
 G=\F_p^{2r},\qquad V=\F_p^r,
\]
and let the maps $F_y:G\to V$ be those of
Theorem~\ref{thm:linear-family}.  Define
\[
 S=\bigcup_{y\in\mathcal Y}\ker F_y\subseteq G.
\]
Let $\Gamma_S$ be the simple graph on $G$ in which distinct $x,x'$ are
adjacent exactly when $x+x'\in S$.

\begin{proposition}[Group palette]\label{prop:group-palette}
With $d,s=o(r)$ as above,
\[
 |S|\le p^{r+d},
 \qquad
 \alpha(\Gamma_S)<p^{r+s}.
\]
\end{proposition}

\begin{proof}
Every retained map is surjective, so every kernel has size $p^r$.
The definition of $S$ therefore gives
\[
 |S|\le|\mathcal Y|p^r\le p^{r+d}.
\]

Suppose that $A\subseteq G$ is independent.  Fix $y\in\mathcal Y$.  If
$z\ne0$ and both $z$ and $-z$ occur in $F_y(A)$, choose $a,a'\in A$ with
$F_y(a)=z$ and $F_y(a')=-z$.  Since $p$ is odd, $z\ne-z$, so $a\ne a'$;
but then $a+a'\in\ker F_y\subseteq S$, a contradiction.  The zero fibre
contains at most one member of $A$, because two distinct members of that
fibre would also have their sum in $\ker F_y$.  The nonzero elements of
$V$ form $(p^r-1)/2$ antipodal pairs.  Hence
\[
 |F_y(A)|\le\frac{p^r+1}{2}.
\]

Any distribution $P$ supported on $T\subseteq V$ satisfies
\[
 \TV(P,U_V)\ge1-\frac{|T|}{|V|},
\]
by testing the event $T$.  If $|A|\ge p^{r+s}$ and $X$ is uniform on
$A$, then $H_{\infty,p}(X)\ge r+s$, while the preceding support bound gives,
for every retained seed,
\[
 \TV(F_y(X),U_V)
 \ge\frac{p^r-1}{2p^r}\ge\frac13.
\]
This contradicts Theorem~\ref{thm:linear-family}.  Therefore
$|A|<p^{r+s}$.
\end{proof}

The zero-fibre argument is where the absence of loops must be handled
carefully: it forbids two distinct zero-fibre points, but it does not
forbid a single point $a$ merely because $2a\in S$.

\section{Carry closure and ordinary integer sums}\label{sec:carries}

Set
\[
 M=p^{2r}
\]
and identify $[M]_0$ with $\F_p^{2r}$ by the base-$p$ digit bijection
\[
 \psi(x)=(x_0,\ldots,x_{2r-1}),
 \qquad x=\sum_{i=0}^{2r-1}x_i p^i.
\]
For $u\in\F_p^{2r}$, define its ordinary carry fibre by
\begin{equation}\label{eq:carry-fibre}
 \calC(u)=\{x+x':0\le x,x'<M,\ \psi(x)+\psi(x')=u\},
\end{equation}
where addition inside $\psi(x)+\psi(x')$ is coordinatewise modulo $p$.

\begin{lemma}[Carry count]\label{lem:carry}
For every $u\in\F_p^{2r}$,
\[
 |\calC(u)|\le2^{2r}.
\]
\end{lemma}

\begin{proof}
In the ordinary addition of $x$ and $x'$, let $c_i\in\{0,1\}$ be the
carry entering digit $i$, with $c_0=0$, and let $c_{2r}$ be the final
carry.  Once $u$ and the vector
\[
 (c_1,\ldots,c_{2r})\in\{0,1\}^{2r}
\]
are fixed, the $i$th ordinary output digit is forced to be
$(u_i+c_i)\bmod p$, and the final digit is $c_{2r}$.  Thus the resulting
integer sum is uniquely determined.  Some carry vectors may be
inconsistent with any pair of summands, which only reduces the count.
\end{proof}

Define the integer palette
\[
 B_M=\bigcup_{u\in S}\calC(u)\subseteq\{0,\ldots,2M-2\}.
\]
Lemma~\ref{lem:carry} and Proposition~\ref{prop:group-palette} give
\[
 |B_M|\le2^{2r}p^{r+d}.
\]
Let $K_{B_M}$ be the simple graph on $[M]_0$ in which distinct integers
$x,x'$ are adjacent when $x+x'\in B_M$.  If $x,x'$ form an edge of
$\Gamma_S$ under the digit bijection, then
$\psi(x)+\psi(x')\in S$, so the carry-fibre definition
\eqref{eq:carry-fibre} and the definition of $B_M$ place their ordinary sum
in $B_M$.  Therefore $K_{B_M}$ contains the group
graph and
\[
 \alpha(K_{B_M})\le p^{r+s}.
\]
No carry has been discarded, and only pairs of distinct vertices are
used in this graph comparison.

\section{Parameters for every interval length}\label{sec:all-N}

We now choose the field and dimension separately for each $N$.  This
also gives the quantitative error term in Theorem~\ref{thm:main}.  Let
$N$ be sufficiently large, put
\[
 L=\log N,
 \qquad
 \lambda_0=\left(\frac{L}{\log L}\right)^{1/3},
 \qquad
 P=\left\lceil\exp(\lambda_0)\right\rceil,
\]
and choose a prime $p\in[P,2P]$, which exists by Bertrand's postulate
\cite{HardyWright2008}.
Set
\begin{equation}\label{eq:r-choice}
 r=\left\lceil\frac{\log N}{2\log p}\right\rceil,
 \qquad M=p^{2r}.
\end{equation}
Then $p$ is odd and tends to infinity, and
\[
 \lambda:=\log p=\lambda_0+O(1),
 \qquad
 r=\Theta(L/\lambda).
\]
Since $\log r=O(\log L)=o(\lambda)$, one has
$\log(pr)=O(\lambda)$, and \eqref{eq:ell-bound} gives
$\ell=O(\lambda)$.  Moreover,
\[
 \bigl(\log(pr)\bigr)^2\log\bigl(2+\log(pr)\bigr)
 =O(\lambda^2\log\lambda)=o(L),
\]
whereas $r\log p=\Theta(L)$.  Thus the hypothesis of
Theorem~\ref{thm:linear-family} holds with room to spare.
The design construction and the quantitative bounds in that theorem give
\[
 \begin{aligned}
 D
 &=O\left(\lambda\sqrt{L/\lambda}
          +\lambda^2\log\lambda\right)
  =O\left(L^{2/3}(\log L)^{1/3}\right),\\
 d,s
 &=O\left(L^{1/3}(\log L)^{2/3}\right),\\
 \frac dr,\frac sr
 &=O\left(\left(\frac{\log L}{L}\right)^{1/3}\right).
 \end{aligned}
\]
Also, by the ceiling in \eqref{eq:r-choice},
\[
 N\le M<Np^2.
\]

Build the $M$-vertex palette $B_M$ above and restrict to the first $N$
vertices.  Delete the normalized sums which cannot come from two distinct
members of $[N]_0$ by putting
\begin{equation}\label{eq:BN}
 B_N=B_M\cap\{1,\ldots,2N-3\}.
\end{equation}
Let $\Gamma_N$ be the subgraph of $\Gamma_S$ induced by $[N]_0$.  Every edge of
$\Gamma_N$ has ordinary sum in
$B_M\cap\{1,\ldots,2N-3\}=B_N$, so
\[
 \Gamma_N\subseteq H_{B_N}.
\]
Therefore
\[
 \alpha(H_{B_N})\le\alpha(\Gamma_N)
 \le\alpha(\Gamma_S)<p^{r+s}.
\]
Together with the palette-size bound, this gives
\[
 |B_N|\le2^{2r}p^{r+d},
 \qquad
 \alpha(H_{B_N})\le p^{r+s}.
\]

We estimate these expressions relative to $N$.  Put $\lambda=\log p$.
By \eqref{eq:r-choice}, $r=L/(2\lambda)+O(1)$.  At these parameters the
bounds for $d$ and $s$ give
$d\lambda,s\lambda=O(D+\log(pr))=O(D+\lambda)$, so
\begin{align*}
 \frac{\log(2^{2r}p^{r+d})}{L}
 &=\frac{2r\log 2}{L}+\frac{(r+d)\lambda}{L}
  \le\frac12+O\left(\frac1\lambda+\frac DL+\frac\lambda L\right),\\
 \frac{\log(p^{r+s})}{L}
 &\le\frac12+O\left(\frac DL+\frac\lambda L\right).
\end{align*}
Here the carry factor contributes
$2r\log2/L=(\log2)/\lambda+O(L^{-1})$, and the ceiling in
\eqref{eq:r-choice} contributes $O(\lambda/L)$.  By the choices above,
\[
 \frac1\lambda,\frac DL,\frac\lambda L
 =O\left(\left(\frac{\log L}{L}\right)^{1/3}\right).
\]
Therefore, for an absolute constant $C_1$ and all sufficiently large $N$,
\begin{equation}\label{eq:BN-final}
 \begin{aligned}
 |B_N|&\le N^{\frac12+C_1
   \left(\frac{\log\log N}{\log N}\right)^{1/3}},\\
 \alpha(H_{B_N})&\le N^{\frac12+C_1
   \left(\frac{\log\log N}{\log N}\right)^{1/3}}.
 \end{aligned}
\end{equation}

\section{Completion of the proof}\label{sec:completion}

Take $N=n-1$ and the normalized palette $B_N$ from
\eqref{eq:BN}.  Select the original sums
\[
 B=\{\,2n+2+a:a\in B_N\,\}.
\]
Because $1\le a\le2N-3=2n-5$, this is a subset of
\[
 \{2n+3,\ldots,4n-3\}\subseteq J_n.
\]
The translation $x\mapsto n+1+x$ identifies the graph generated by $B_N$
on $[N]_0$ with $G_B$ on $I_n$.  Hence \eqref{eq:BN-final} gives
\[
 |B|+\alpha(G_B)
 \le n^{\frac12+O\!\left(
   \left(\frac{\log\log n}{\log n}\right)^{1/3}\right)}.
\]
The factor $2$ from adding the two bounds contributes only
\[
 \frac{\log2}{\log n}
 =o\left(\left(\frac{\log\log n}{\log n}\right)^{1/3}\right)
\]
to the exponent.  Replacing the scale $N=n-1$ by $n$ changes the exponent
by a smaller amount.  These errors are absorbed in the displayed term.
By Proposition~\ref{prop:minmax}, this is the desired upper bound for
$f(n)$.  Proposition~\ref{prop:loglower} supplies the lower bound, and
Theorem~\ref{thm:main} follows.

For clarity about the quantifiers, the construction above chooses one
set $B$ for each $n$.  The preceding estimate for
$|B|+\alpha(G_B)$ therefore holds with $\alpha(G_B)$ replaced by
$|C|$ for \emph{every} $B$-admissible $C$.  Since the entire exponent
correction displayed above tends to zero, it is smaller than $\varepsilon$
for sufficiently large $n$.  This proves the stated assertion for every
fixed $\varepsilon>0$ and every sufficiently large $n$, not merely along a
subsequence.

\section{Concluding remarks}

The proof separates three interfaces that are easy to conflate.  The
extractor has $p$-ary linear outputs but binary code coordinates and a
binary seed; the weak-design overlap moment is therefore $2^t$, not $p^t$.
Ordering the design sets exposes a further resource: at hybrid row $i$, the
remaining $r-i$ output coordinates pay for the excess above the unavoidable
baseline overlap.
The union of the resulting kernels controls finite-field sum graphs, and
the carry closure in Section~\ref{sec:carries} then transfers that control
to ordinary integer addition.  Keeping these interfaces explicit yields
the quantitative exponent correction
$O((\log\log n/\log n)^{1/3})$.

The argument determines the exponent of $f(n)$ but not its finer order of
growth.  The universal lower bound is $\Omega(\sqrt{n\log n})$, whereas the
constructed upper bound is
\[
 \sqrt n\,\exp\!\bigl(
 O((\log n)^{2/3}(\log\log n)^{1/3})\bigr).
\]
Closing this remaining subpolynomial gap would require a sharper palette,
a stronger universal lower bound, or both.

\begingroup
\small
\begin{thebibliography}{99}

\bibitem{ErdosProblems788}
T.~F. Bloom,
\newblock Erd\H{o}s Problem \#788,
\newblock \url{https://www.erdosproblems.com/788},
\newblock accessed July 18, 2026.

\bibitem{Erdos1973}
P.~Erd\H{o}s,
\newblock Problems and results on combinatorial number theory,
\newblock in \emph{A Survey of Combinatorial Theory}
(J.~N. Srivastava et al., eds.; Proc. Internat. Sympos.,
Colorado State Univ., Fort Collins, 1971),
North-Holland, 1973, pp.~117--138,
\newblock \href{https://doi.org/10.1016/B978-0-7204-2262-7.50017-X}
{doi:10.1016/B978-0-7204-2262-7.50017-X}.

\bibitem{Choi1971}
S.~L.~G. Choi,
\newblock On a combinatorial problem in number theory,
\newblock \emph{Proc. London Math. Soc.} (3) \textbf{23} (1971), 629--642,
\newblock \href{https://doi.org/10.1112/plms/s3-23.4.629}
{doi:10.1112/plms/s3-23.4.629}.

\bibitem{BSS2000}
A.~Baltz, T.~Schoen, and A.~Srivastav,
\newblock Probabilistic construction of small strongly sum-free sets via
large Sidon sets,
\newblock \emph{Colloq. Math.} \textbf{86}, no.~2 (2000), 171--176,
\newblock \href{https://doi.org/10.4064/cm-86-2-171-176}
{doi:10.4064/cm-86-2-171-176}.

\bibitem{AlonPham2025}
N.~Alon and H.~T. Pham,
\newblock Random Cayley graphs and random sumsets,
\newblock arXiv:2509.02561, 2025,
\newblock \url{https://arxiv.org/abs/2509.02561}.

\bibitem{AKS1999}
N.~Alon, M.~Krivelevich, and B.~Sudakov,
\newblock Coloring graphs with sparse neighborhoods,
\newblock \emph{J. Combin. Theory Ser. B} \textbf{77} (1999), 73--82,
\newblock \href{https://doi.org/10.1006/jctb.1999.1910}
{doi:10.1006/jctb.1999.1910}.

\bibitem{RRV2002}
R.~Raz, O.~Reingold, and S.~Vadhan,
\newblock Extracting all the randomness and reducing the error in
Trevisan's extractors,
\newblock \emph{J. Comput. System Sci.} \textbf{65} (2002), 97--128,
\newblock \href{https://doi.org/10.1006/jcss.2002.1824}
{doi:10.1006/jcss.2002.1824}.

\bibitem{Trevisan2001}
L.~Trevisan,
\newblock Extractors and pseudorandom generators,
\newblock \emph{J. ACM} \textbf{48} (2001), 860--879,
\newblock \href{https://doi.org/10.1145/502090.502099}
{doi:10.1145/\allowbreak 502090.502099}.

\bibitem{CoverThomas2006}
T.~M. Cover and J.~A. Thomas,
\newblock \emph{Elements of Information Theory}, 2nd ed.,
Wiley-Interscience, Hoboken, NJ, 2006.

\bibitem{HardyWright2008}
G.~H. Hardy and E.~M. Wright,
\newblock \emph{An Introduction to the Theory of Numbers}, 6th ed.,
revised by D.~R. Heath-Brown and J.~H. Silverman,
Oxford University Press, Oxford, 2008.

\end{thebibliography}
\endgroup

\end{document}
